\documentclass[11pt,a4paper]{article}
\usepackage{DDTA_METHODOLOGY_GUIDE_STYLE_R1}

\providecommand{\candidatetag}{\ddtatag{DDTAAmber}{CANDIDATE / NOT ADMITTED}}
\providecommand{\rejectedtag}{\ddtatag{DDTARed}{R25 REJECTED}}

\begin{document}
\selectlanguage{italian}
\hypersetup{pdftitle={DDTA Documentation Authoring Guide - Revision 6 Candidate R2}}
\fancyhead[L]{\sffamily\small\color{DDTAGray}DDTA Documentation Authoring Guide}
\fancyhead[R]{\sffamily\small\color{DDTAAmber}R6 CANDIDATE R2 - R43}
\fancyfoot[L]{\sffamily\scriptsize\color{DDTAGray}R5 remains current authority; this document is non-normative}

\begin{titlepage}
\thispagestyle{empty}
\vspace*{8mm}
{\sffamily\bfseries\color{DDTANavy}\fontsize{15}{18}\selectfont Documentation-Driven Threat Analysis}\par
\vspace{5mm}
{\sffamily\bfseries\color{DDTANavy}\fontsize{29}{33}\selectfont Documentation Authoring Guide}\par
\vspace{2mm}
{\sffamily\color{DDTABlue}\fontsize{18}{22}\selectfont Revision 6 - Candidate documentation-only consolidation R2}\par
\vspace{9mm}
\begin{tcolorbox}[enhanced,arc=3pt,boxrule=0pt,colback=DDTANavy,left=12pt,right=12pt,top=11pt,bottom=11pt]
{\color{white}\sffamily\bfseries\large Authoring DDTA nativo: project documentation prima dei consumer downstream}\par
\vspace{2mm}
{\color{white!88}\small Consolidamento documentation-only della R5: preservazione informativa, source gaps, condizioni, selection semantics, pipeline/order e separazione rigorosa tra project documentation e consumer downstream. La Base Analysis compare soltanto come handoff/feedback boundary.}
\end{tcolorbox}
\vspace{8mm}
\begin{tabularx}{\textwidth}{L{45mm}Y}
\toprule
\textbf{Stato} & \texttt{CANDIDATE CUMULATIVE / NON-NORMATIVE / REVIEW ONLY} \\
\textbf{Authority corrente preservata} & \nolinkurl{DDTA_DOCUMENTATION_BA_AUTHORING_GUIDE_R5.tex} \\
\textbf{Scopo} & L1 conceptual meaning -> L2 readable governed documentation; L3/L4 come consumer, non authority. \\
\textbf{Repository baseline R43} & \texttt{bea1e6567d3f861294378b201648b0853b3c8a43} \\
\textbf{Data} & 12 settembre 2026 \\
\bottomrule
\end{tabularx}
\vfill
\begin{ddtaopen}[title={Authority stop}]
La R5 resta current methodology authority. Questa R6 candidate riguarda esclusivamente la metodologia di documentazione: non definisce operatori, signature o cardinalita' della Base Analysis e non introduce automaticamente nuovi tipi L1, campi o project requirement. La promotion richiede review separata.
\end{ddtaopen}
\end{titlepage}

\tableofcontents
\clearpage

\section{Scopo, lineage e come usare questa guida}

La predecessor authority conserva il filename storico \nolinkurl{DDTA_DOCUMENTATION_BA_AUTHORING_GUIDE_R5.tex}. In R6 la denominazione viene corretta in \textbf{Documentation Authoring Guide}: il cambio di nome chiarisce il confine metodologico e non modifica retroattivamente R5.

L'obiettivo resta quello della R5: spiegare il \textbf{significato concettuale L1} dei modelli DDTA e delle loro relazioni affinche' un autore possa produrre una \textbf{rappresentazione L2} fedele, leggibile e governabile.

\begin{ddtacore}[title={Cumulative preservation}]
Le regole chiuse e qualificate della R5 restano preservate. R43 preserva i learning documentali emersi dal ciclo R25 e li separa formalmente dal metodo BA. Nessuna lesson learned viene promossa a contratto solo perche' e' utile o ricorrente.
\end{ddtacore}

\begin{tabularx}{\textwidth}{L{60mm}Y}
\toprule
\textbf{Etichetta} & \textbf{Significato} \\
\midrule
\closedtag & R5/precedessor rule mantenuta. \\
\holdouttag & Evidence R25 che rafforza una pratica senza cambiare da sola L1. \\
\candidatetag & Candidate downstream construct non ammesso; puo' essere citato solo come boundary di non-authority e non deve contaminare la project documentation. \\
\rejectedtag & Working hypothesis gia' respinta; non va mantenuta come requisito futuro. \\
\opentag & Pressure reale ancora da chiudere; deve restare visibile ma non essere inventata nella documentazione. \\
\layertwotag & Convenzione di rappresentazione leggibile, non nuova semantica L1. \\
\bottomrule
\end{tabularx}

\section{Perimetro metodologico: Documentation Authoring only}

\begin{ddtacore}[title={Separation contract}]
Questa guida governa esclusivamente come riconoscere, strutturare, revisionare e promuovere la \textbf{project documentation DDTA}. Non e' una guida alla Base Analysis e non duplica il relativo operator vocabulary. La Base Analysis entra qui soltanto in due punti: \textbf{handoff} della documentazione accettata verso un consumer downstream e \textbf{feedback} downstream che puo' proporre una candidate documentation change, soggetta a nuova governance.
\end{ddtacore}

\begin{tabularx}{\textwidth}{L{47mm}Y}
\toprule
\textbf{Dentro questa guida} & \textbf{Fuori da questa guida} \\
\midrule
MacroRequirement, Decision, Requirement/FR, SpecializedRequirement, SecurityRequirement & BA operator definitions/signatures/roleKey/cardinalities \\
GovernedDocument, authority, lifecycle, provenance & BAReferent / BAProposition operational construction \\
L1/L2 authoring, L3/L4 boundary, STOP rules & operator selection algorithms and BA projection mechanics \\
source gaps, ambiguity preservation, semantic-owner review & threat-analysis method, DFD topology and STRIDE mapping \\
conditions/selection/pipeline meaning when project-governed & candidate BA constructs as documentation vocabulary \\
\bottomrule
\end{tabularx}

\begin{ddtaanti}[title={No duplicated BA method}]
Se una domanda riguarda ``quale operatore usare'', ``quale roleKey assegnare'', ``quale signature applicare'' o ``come costruire una BAProposition'', il lettore deve passare alle guide BA Core/Complete. Questa guida deve solo assicurare che il significato documentale necessario sia governato, leggibile e non inventato.
\end{ddtaanti}

\section{Native DDTA authoring e validazione di documentazione esistente}
\closedtag

Il processo normale parte dal problema e dal significato governato del progetto.

\begin{lstlisting}
NORMAL DDTA AUTHORING

project problem / governed meaning
        |
        v
authority and governance
        |
        v
MacroRequirement -> Decision -> FunctionalRequirement
                                   |
                                   v
                         SpecializedRequirement
                                   |
                                   v
                          SecurityRequirement
\end{lstlisting}

La ricostruzione da documentazione preesistente e' invece un protocollo di validazione: recupera solo significato supportato, applica le normali domande DDTA e misura preservazione, chiarezza, gap e representability.

\begin{ddtaanti}[title={Authority inversion}]
Non modificare la documentazione perche' BA, STRIDE, validator, LLM o tool ``hanno bisogno'' di un elemento. Un consumer downstream puo' esporre una domanda o un gap; non puo' inventare la risposta di progetto.
\end{ddtaanti}

\section{Layering: L1, L2, L3 e L4}
\closedtag

\begin{tabularx}{\textwidth}{L{14mm}L{42mm}Y}
\toprule
\textbf{Layer} & \textbf{Responsabilita'} & \textbf{Esempio} \\
\midrule
L1 & concetti e semantica del metamodel & un FR ha una parent Decision; SecurityRequirement IS-A SpecializedRequirement. \\
L2 & rappresentazione canonica leggibile & file Markdown/LaTeX/YAML, Context, normative clauses, parent references. \\
L3 & realizzazione del modello & schema, indici, validator, serializzazione e projection. \\
L4 & tool support e conformance & editor, suggestion, validation UX, conformance declaration. \\
\bottomrule
\end{tabularx}

\begin{ddtacore}[title={Direzione}]
L1 governa il significato; L2 lo rende leggibile; L3 lo realizza; L4 lo supporta. Un bisogno di L3/L4 non dimostra una nuova proprieta' L1.
\end{ddtacore}

\section{GovernedDocument e authority}
\closedtag

Un \texttt{GovernedDocument} e' un elemento documentale il cui significato e' riconosciuto dal processo di governance: ha identita' distinguibile nel tempo, lifecycle interpretabile e una collocazione rispetto all'authority corrente.

\begin{lstlisting}
source file / code / config / test / model / dataset
                    !=
              GovernedDocument
\end{lstlisting}

Un artefatto tecnico puo' essere evidence reale senza diventare Decision o Requirement. Diventa project commitment solo quando la governance gli attribuisce quel significato.

\subsection{Identity, lifecycle, authority e provenance}

\begin{lstlisting}
ID != hierarchy
ID != chronology
ID = stable referent

Lifecycle = candidate / current / superseded / retired / ...
Gate      = PASS / REWORK / STOP / REOPEN / ...
\end{lstlisting}

Recency non e' authority. Provenance spiega perche' un significato e' stato introdotto o modificato; non sostituisce il contenuto normativo.

\begin{ddtaexample}[title={DermaTriage FR-11}]
FR-11 e' una identita' storica superseded. Lo split in piu' obbligazioni indipendenti non autorizza a riciclare silenziosamente FR-11 per una sola parte: l'identita' storica resta preservata e i prodotti dello split ricevono identita' proprie.
\end{ddtaexample}

\section{Principi trasversali}

\subsection{Authority prima dell'authoring}
\closedtag
Prima di creare o revisionare un elemento DDTA: identificare source set e baseline autorizzati; distinguere current/candidate/working/historical/superseded; preservare provenance; non usare recency come authority; non permettere a BA o analysis di diventare source of project truth.

\subsection{Documentazione per lettori di progetto}
\layertwotag
La documentazione serve prima a autori e reviewer del progetto. Evitare mechanics BA o threat-analysis nel corpo se non sono significato del progetto.

\subsection{Minimum sufficient governed meaning}
\closedtag
Documentare la minima conoscenza governata sufficiente a stabilizzare il significato, rendere verificabile la responsabilita' e permettere una successiva normalizzazione senza inventare dettagli.

\begin{ddtacheck}[title={Stopping criterion}]
Abbiamo abbastanza semantica per preservare la distinzione materiale o rispondere alla domanda dichiarata nello scope corrente? YES: stop sul branch. NO: cercare la smallest missing distinction; non aggiungere dettaglio per completezza narrativa.
\end{ddtacheck}

\section{Unified DDTA native authoring flow}

\begin{lstlisting}
AUTHORITY GATE
    |
PROJECT PROBLEM FRAMING
    |
MacroRequirement
    |-- sufficient? YES -> STOP AT MR
    |
Decision
    |-- no further operational obligation? -> STOP + REVIEW
    |
FunctionalRequirement
    |-- no independent additional property? -> STOP AT FR
    |
SpecializedRequirement
    |-- security property? YES -> SecurityRequirement
    `-- NO -> preserve specialization pressure; do not invent subtype
\end{lstlisting}

\section{Step 0 - Project problem framing}
\closedtag

Il framing e' una precondizione di metodo, non un nuovo tipo L1.

\begin{ddtacheck}[title={Domanda fondamentale}]
Quale problema o classe di problemi deve affrontare il progetto, entro quale boundary significativo, rimuovendo temporaneamente le scelte di soluzione correnti?
\end{ddtacheck}

\begin{ddtaexample}[title={DermaTriage}]
Supportare il triage precoce di casi dermatologici relativi a lesioni potenzialmente oncologiche, usando le informazioni disponibili per determinare priorita' e supportare l'instradamento specialistico. Il framing non richiede di nominare EfficientNet, Qwen, ChromaDB o BioMistral.
\end{ddtaexample}

\section{Step 1 - MacroRequirement}
\closedtag

Un \MR{} governa una sola responsabilita' macro durevole necessaria ad affrontare il problema.

Campi L2 da interpretare semanticamente: Title, Intent, Context, Stakeholders, Scope, Assumptions/Constraints e \texttt{dependsOn}. \texttt{dependsOn} e' non-gerarchico e non deve esprimere semplice ordine narrativo.

\begin{ddtaexample}[title={DermaTriage MR-02 - STOP AT MR}]
La source governa la responsabilita' macro di specialist-oriented routing ma non stabilisce ancora vocabolario destinazione, selection rule, input, fallback o booking responsibility. Il risultato corretto e' STOP AT MR: non si inventano Decision o FR per riempire la gerarchia.
\end{ddtaexample}

\section{Step 2 - Semantic-sufficiency review}
\closedtag

Il gate non chiede se il testo potrebbe essere piu' dettagliato; chiede se possono sopravvivere letture materialmente diverse.

Evidence classification:
\begin{lstlisting}
AFFIRMED
DENIED
NOT SPECIFIED
CONFLICTING
\end{lstlisting}

\NS{} non equivale a FALSE. \texttt{CONFLICTING} non autorizza a scegliere la lettura piu' comoda.

\begin{ddtaexample}[title={DermaTriage deployment authority}]
FR-09 qualifica un candidate adaptation rispetto a criteri governati, ma non stabilisce automaticamente chi autorizzi il deployment finale. ``Qualification autorizza deployment automatico'' resta \texttt{NOT SPECIFIED}, non viene impostato a TRUE o FALSE per comodita'.
\end{ddtaexample}

\section{Step 3 - Decision}
\closedtag

Una \DEC{} e' un commitment progettuale significativo che restringe esattamente un MR fissando posizione, boundary, policy, convention, strategy, technology/architecture choice o altro commitment materialmente rilevante.

\begin{ddtacore}[title={Consequences descrive il sistema}]
\texttt{Consequences} documenta effetti, trade-off, responsabilita' o obblighi downstream. Non deve contenere istruzioni metodologiche come ``rieseguire il gate'', ``revalidare'' o ``la BA dovra'...''.
\end{ddtacore}

\subsection{Technical choice routing}
\holdouttag
Se il progetto governa materialmente una scelta concreta, il nome di componente, sistema, software, modello, algoritmo o protocollo deve restare presso il semantic owner corretto. La neutralizzazione temporanea serve a classificare, non a genericizzare la documentazione finale.

\begin{ddtaexample}[title={DermaTriage DEC-12}]
L'MR resta generale, mentre DEC-12 governa la pipeline AI image-based a quattro stadi. Gli FR sotto la Decision conservano i referent concreti EfficientNet-B4, Qwen2-VL-7B-Instruct, ChromaDB/all-MiniLM-L6-v2 e BioMistral-7B quando partecipano materialmente al comportamento governato.
\end{ddtaexample}

\section{Step 4 - FunctionalRequirement}
\closedtag

Un \FR{} e' un obbligo operativo governato e independently assessable che discende da esattamente una Decision.

\begin{lstlisting}
Requirement [abstract]
    normativeClause : NormativeClause [1..*]

FunctionalRequirement
    extends Requirement
    parentDecision : Decision [1]
\end{lstlisting}

La normative prose resta primary source. SPO reference e label L2 aiutano lettura e tooling ma non possono perdere condition, lifecycle, correlation, failure o result semantics.

\begin{ddtacore}[title={Specifico non significa non-funzionale}]
Un FR resta funzionale anche se nomina un componente, modello, algoritmo, sistema o interfaccia concreto. Se il referent e' governato e materialmente parte del comportamento, genericizzarlo puo' perdere informazione.
\end{ddtacore}

\subsection{Allowed result domain vs conditional selection rule}
\closedtag
Se la source governa solo outcome ammessi, non inventare il mapping che seleziona l'outcome. Se governa un vero mapping condizionale, quel mapping appartiene alla semantica dell'FR.

\begin{ddtaexample}[title={DermaTriage FR-02}]
\begin{lstlisting}
HIGH + confidence > 0.85 -> P1
HIGH                     -> P2
MEDIUM                   -> P3
LOW                      -> P4
\end{lstlisting}
Questa e' una vera regola condizionale governata, non una lista di outcome.
\end{ddtaexample}

\section{Step 5 - Requirement, normativeClause e split}
\closedtag

\begin{lstlisting}
1 Requirement != 1 textual sentence
1 Requirement  = 1 coherent normative obligation
\end{lstlisting}

Split quando clausole possono essere introdotte, revisionate, ritirate o assessed indipendentemente senza cambiare l'identita' normativa delle rimanenti. Il numero di righe o test case non decide lo split.

\section{Step 6 - SpecializedRequirement}
\closedtag

\SR{} e' una obbligazione concern-specific aggiuntiva che vale insieme al parent FR senza ridefinire la funzione ordinaria o scegliere la realization.

\begin{lstlisting}
SpecializedRequirement [abstract]
    extends Requirement
    parentFunctionalRequirement : FunctionalRequirement [1]
\end{lstlisting}

Applicare removal test e autonomous capability test. Non inventare concrete subtype solo per classificare un caso osservato.

\section{Step 7 - SecurityRequirement}
\closedtag

\SECR{} e' uno \SR{} con una singola security property governante e failure mode identificabile nelle clauses.

\begin{lstlisting}
SecurityRequirement
    extends SpecializedRequirement
    protectedSecurityProperty : SecurityProperty [1]
\end{lstlisting}

Property, mechanism e attack cause restano distinti. TLS, signature o certificate non sono implicati automaticamente da Confidentiality, Integrity o provenance.

\section{Step 8 - Cross-MR e consumed-service boundary}
\closedtag

\begin{lstlisting}
consumption != ownership
consumption != second parent
shared service != shared FR ownership
\end{lstlisting}

Dependency non significa containment o co-ownership. Una property richiesta dal consumer puo' essere assunta dal service solo se evidence governata la garantisce.

\section{Step 9 - Terminologia, information binding e tool authority}
\closedtag

Prosa naturale e canonical referent possono convivere. Nomi propri, processi, algoritmi, sistemi, protocolli, dati e binding materialmente governati non devono essere cancellati durante l'astrazione.

\begin{ddtacheck}[title={Information-preservation check}]
Un elemento materialmente governato e' stato genericizzato o eliminato per far apparire la documentazione piu' astratta? Se SI, riportarlo presso il semantic owner corretto. La normalizzazione dei sinonimi e' compito della BA, non dell'authoring mediante perdita di informazione.
\end{ddtacheck}

\section{Step 10 - Downstream semantic propagation}
\closedtag

Una correzione upstream richiede review non solo per contraddizione diretta ma anche per silent precondition, lost governed distinction, semantic narrowing/widening e misplaced ownership.

\section{Step 11 - Diagnostics: source gap, guide problem o metamodel pressure}
\holdouttag

\begin{longtable}{L{47mm}L{93mm}}
\toprule
\textbf{Diagnosi} & \textbf{Interpretazione} \\
\midrule\endhead
Source/documentation gap & significato necessario non governato, ambiguo o conflittuale nella source. \\
Authoring-guide problem & significato rappresentabile ma domande/gate non guidano stabilmente l'autore. \\
Metamodel pressure & significato chiaro e governato ma costrutti correnti non lo rappresentano onestamente. \\
Representation/L2 issue & L1 sufficiente ma forma canonica leggibile/pratica incompleta. \\
Downstream/BA issue & documentazione sufficiente ma rappresentazione analitica perde significato. \\
No gap & differenza presentazionale o dettaglio fuori scope. \\
\bottomrule
\end{longtable}

\begin{ddtacore}[title={Minimum owning layer}]
Prima classificare il problema, poi riaprire solo il minimo layer che semanticamente possiede il significato. Una improvement utile non e' automaticamente difetto del metamodel.
\end{ddtacore}

\section{Step 12 - Technical, configuration e verification evidence}
\holdouttag

La preservazione dell'informazione tecnica non introduce un technical-document model parallelo. Il fatto materialmente governato va nei normali MR/Decision/Requirement; il fatto non governato resta evidence/configuration/gap.

\subsection{Parameter governance boundary}
\holdouttag
\begin{lstlisting}
source numeric/config value
    -> semantic owner
    -> lifecycle / purpose
    -> requirement/property classification
    -> parameter boundary
    -> current concrete binding
\end{lstlisting}

Same literal non significa same parameter: owner, lifecycle e purpose partecipano all'identita' semantica.

\section{Decision-to-Requirement completeness test}
\holdouttag

\begin{ddtacheck}[title={Counterexample gate}]
Assumiamo che tutti i Requirement downstream correnti siano soddisfatti. La parent Decision puo' comunque essere violata? YES: downstream coverage incompleta; cercare obligation mancante o specialization non rappresentata. NO: coverage plausibilmente sufficiente, subject to altri gate.
\end{ddtacheck}

Il test non autorizza a creare Requirement artificiali quando la source non governa il significato mancante.

\section{Step 13 - Documentation completeness e promotion gate}
\closedtag

Prima della promotion verificare branch coerenti, propagation, \NS{} preservati, assenza di research commentary nella baseline, semantic owner stabiliti, nomi/distinzioni governate non persi, gap bloccanti risolti o esclusi esplicitamente, governance di promotion e authority registry aggiornato.

\section{Step 14 - Handoff ai consumer downstream}
\closedtag

La project documentation resta authority. L'handoff verso Base Analysis, security analysis, test design o altri consumer deve offrire authority/baseline, source set, termini/referent governati, relazioni/condizioni materialmente necessarie, diagnostics/\NS{} e revalidation context. Il consumer sceglie poi la propria rappresentazione senza riscrivere la source.

\begin{ddtacheck}[title={R43 handoff target}]
Per la Base Analysis, il consumer metodologico e' ora rappresentato dal candidate guide set separato: \texttt{DDTA\_BASE\_ANALYSIS\_CORE\_GUIDE} per semantica/sintassi e \texttt{DDTA\_BASE\_ANALYSIS\_COMPLETE\_GUIDE} per domande ed esempi. Questa reference stabilisce un confine; non promuove quelle candidate a current authority.
\end{ddtacheck}

\begin{ddtaprinciple}[title={View generation}]
DFD, diagrammi architetturali e altre viste non diventano automaticamente documentation authority. Quando servono, e' preferibile materializzarli deterministicamente dalla BA o da projection trace-bound, evitando una seconda source manuale di significato.
\end{ddtaprinciple}

\section{Step 15 - Downstream feedback senza authority inversion}
\closedtag

\begin{lstlisting}
accepted documentation
    -> Base Analysis
    -> security analysis / Finding
    -> candidate documentation change
    -> governed review
    -> if accepted: updated GovernedDocument
    -> rebuilt/revalidated Base Analysis
\end{lstlisting}

Finding non equivale a Requirement. Threat analysis non canonizza automaticamente mechanism o obligation.

\section{R25 downstream findings: implicazioni per la qualita' documentale}

Questa sezione non cambia il metamodel: rende visibile cio' che il control cycle ha insegnato all'authoring/review.

\subsection{1. Un operator audit downstream non e' un authoring target}
\holdouttag
L'autore non deve scrivere documentazione per ``coprire'' i 14 operatori. Il risultato R25 e' che l'intera base a 14 operatori e' stata auditata sul corpus; nessun 15esimo operatore e' stato ammesso. Questo e' evidence della BA, non un requisito sulla prosa del progetto.

\subsection{2. Condition semantics devono essere governate}
\holdouttag
Conjunction, comparison e condition satisfaction possono diventare project meaning quando la source le governa. Non introdurre OR, NOT, comparator, threshold o property model soltanto perche' il downstream li sa rappresentare.

\begin{ddtaopen}[title={Negative non-sufficiency}]
``A da solo non e' sufficiente per B'' non equivale a ``NOT B''. Se la source usa linguaggio di non-sufficiency, preservare la frase governata e non trasformarla in una negazione logica piu' forte.
\end{ddtaopen}

\subsection{3. Selection semantics richiedono precisione documentale}
\holdouttag
Quando il progetto governa una scelta tra candidati, distinguere cio' che e' realmente deciso:
\begin{lstlisting}
candidate population
criterion / selection basis
ordering or ranking        [only if governed]
bound                      [only if governed]
selected membership/result
later use/retrieval
\end{lstlisting}

\begin{ddtaexample}[title={DermaTriage FR-18 / FR-06 / FR-20}]
FR-18 governa similarity ordering, bound 5 e selected historical cases; FR-06 governa pertinence, recency, bound 20 e selected evidence set ma lascia dedup/underfill/reuse/overlap non specificati; FR-20 governa il best checkpoint per validation Macro F1 senza richiedere un full reusable ranking.
\end{ddtaexample}

\begin{ddtaanti}[title={Top-K shorthand}]
Non usare \texttt{topK} come se fosse una primitive unica che nasconde candidate population, criterion, ordering, bound e selected result. La documentazione deve dire solo le componenti che il progetto governa.
\end{ddtaanti}

\subsection{4. Pipeline membership/order non e' dependency}
\opentag
Quando il progetto governa una pipeline, documentare esplicitamente cio' che e' vero senza forzare una relation downstream prematura.

\begin{lstlisting}
pipeline membership
    != ordering
    != precedes
    != dependOn
    != transfer
\end{lstlisting}

\begin{ddtaexample}[title={DermaTriage DEC-12}]
La Decision governa una pipeline sequenziale a quattro stadi. E' legittimo documentare membership e sequenza a livello di Decision/FR perche' sono project meaning. Non e' legittimo trasformare automaticamente ogni coppia di stadi in \texttt{dependOn} se la source non governa prerequisite semantics.
\end{ddtaexample}

\subsection{5. Candidate construct downstream non diventano termini della documentazione}
\candidatetag
R25 ha testato positivamente \texttt{provideService}, \texttt{storedIn} e \texttt{initiate}, ma restano non ammessi. L'autore documenta provider/service, at-rest association o initiation quando il progetto li governa, usando linguaggio di dominio. Non deve modellare la prosa per anticipare il nome del candidate operator.

\section{DermaTriage explicit documentation gaps: come preservarli}
\holdouttag

Il corpus corrente espone gap che devono restare project-documentation issue finche' nuova authority non li risolve. Tra questi:
\begin{itemize}
\item diagnostic authority boundary;
\item no-image required/optional symptom inputs e missing-value semantics;
\item P-scale complete input semantics e no-image binding;
\item SLA trigger/outcome/owner/normative binding;
\item predicted pathology obligation/authority;
\item specialist selection vocabulary/rule/input/fallback/booking;
\item review content, lifecycle e disagreement binding;
\item adaptation counting e threshold authority;
\item prompt evidence window ordering/membership/underfill/dedup/reuse/overlap;
\item supervision input invalid/missing/conflicting semantics;
\item deployment authority/automaticity;
\item acceptance e rollback binding;
\item evaluation consistency tra sensitivity/recall/correct-count.
\end{itemize}

\begin{ddtaanti}[title={Gap-to-BA shortcut vietata}]
Non scegliere una rappresentazione downstream e poi riscrivere la documentation per farla combaciare. Il gap deve essere risolto dal semantic owner del progetto o restare esplicitamente gap.
\end{ddtaanti}

\section{Workflow di review A0-A12 preservato}
\closedtag

\begin{longtable}{L{17mm}L{43mm}L{80mm}}
\toprule
\textbf{Gate} & \textbf{Focus} & \textbf{Domanda essenziale} \\
\midrule\endhead
A0 & Authority/source & quali fonti/baseline possono determinare project meaning? \\
A1 & Problem framing & problema comprensibile senza conoscere la soluzione corrente? \\
A2 & MR review & macro-responsabilita' stabili, distinte e sufficienti? \\
A3 & Semantic sufficiency & letture materialmente diverse? proposition minima? evidence state? \\
A4 & Decision review & commitment e responsibility boundary espliciti e non ridondanti? \\
A5 & FR / D3 & ogni Decision operazionalizzata con FR sufficienti e assessable? \\
A6 & Requirement split & ogni Requirement e' una coherent obligation? \\
A7 & SpecializedRequirement & property aggiuntive autonome rispetto al parent FR? \\
A8 & SecurityRequirement & quali specialization sono security property governate? \\
A9 & Cross-MR & consumption, dependency e ownership distinti? \\
A10 & Terminology / bindings & informazioni e current realization preservati senza diventare nuova authority? \\
A11 & Propagation / regression & correzioni upstream propagate senza loss/narrowing/widening? \\
A12 & Completeness / promotion & candidate pronta per promotion e downstream handoff? \\
\bottomrule
\end{longtable}

\section{Checklist cumulativa di authoring}

\begin{longtable}{L{39mm}L{101mm}}
\toprule
\textbf{Area} & \textbf{Domande da non saltare} \\
\midrule\endhead
Governance & identity, lifecycle, authority e provenance sono distinti? Recency viene confusa con authority? \\
Layering & sto descrivendo L1 o una convenzione L2/L3/L4? \\
Framing & problema e boundary restano comprensibili rimuovendo la soluzione corrente? \\
MR & responsabilita' macro stabile? Context elimina ambiguita' materiale? STOP possibile? \\
Decision & vero commitment? parent MR unica? Consequences di sistema, non di metodo? \\
FR & parent Decision unica? coherent operational obligation? normative prose primaria? \\
Split & failure/change/review indipendenti giustificano separate identity? \\
SR/SecR & strengthening autonomo? property governante singola? mechanism separato? \\
Cross-MR & dependency/consumption/ownership corretti? \\
Technical facts & referent concreto e' project commitment o solo realization/evidence? \\
Parameters & semantic owner, lifecycle e purpose chiari? \\
Conditions & comparator/conjunction/negation realmente governati? \\
Selection & candidate population, basis, bound/order/result distinti e source-supported? \\
Pipeline & membership/order/precedes/dependency non collassati? \\
Gaps & source gap registrato invece di completato dal downstream? \\
Completeness & Requirement possono passare mentre parent Decision resta violata? \\
Promotion & governance autorizza davvero la baseline? \\
Downstream handoff & il consumer puo' derivare la propria rappresentazione senza inventare project meaning? \\
\bottomrule
\end{longtable}

\section{Cosa non deve fare la documentazione DDTA}

La documentation non deve:
\begin{itemize}
\item anticipare STRIDE/STRIDE-AI o un altro threat method;
\item scegliere operatori BA o DFD topology per soddisfare un consumer;
\item trasformare test, config, runtime state o model artifact in project truth automaticamente;
\item creare Decision vuote per riempire la gerarchia;
\item inventare threshold, branch, fallback o algoritmo per rendere un FR piu' testabile;
\item trasformare documentation gap in Requirement senza semantic-owner review;
\item usare un field L2 come prova di una proprieta' L1;
\item aggiungere pseudo-campi tecnici paralleli quando il significato e' esprimibile dai normali elementi DDTA;
\item inserire nel documento di progetto istruzioni su review/revalidation/gate/BA;
\item far diventare BA, LLM, tool o threat analysis authority;
\item scrivere in funzione dei candidate operator CC-02/03/04;
\item trasformare pipeline order in dependency senza source support;
\item completare selection semantics mancanti perche' una struttura R25 le prevede.
\end{itemize}

\section{Downstream pressure watchlist: non-authoring contract}
\opentag

Le pressure seguenti provengono dalla review BA R25. Sono riportate solo per evitare che un autore le ``risolva'' inventando project meaning. \textbf{Non fanno parte del contratto di authoring}, non sono campi L1/L2 e la loro soluzione appartiene al metodo BA o al semantic owner del progetto, a seconda del caso:
\begin{itemize}
\item PR-01 function/process/behavior identity / \texttt{perform?};
\item PR-02 pipeline composition/order;
\item PR-04 boundary/interaction association;
\item PR-05/06 ordered/scalar comparison;
\item PR-07 structured information contract;
\item PR-09 acquisition/refresh;
\item PR-12 negative implication/non-sufficiency;
\item PR-14 applicability/configuration binding;
\item CMD-OP04 observe.result removal;
\item CMD-OP05 transition state/value refinement;
\item OBS-OT-01 operation-target/effect-scope.
\end{itemize}

\begin{ddtacore}[title={Authoring response alle open pressure}]
L'autore continua a documentare il significato del progetto in linguaggio DDTA normale. Se una pressure rende impossibile preservare onestamente un significato governato, si registra il counterexample e si riapre il minimo owning layer metodologico; non si distorce il progetto per evitare la pressure.
\end{ddtacore}

\appendix
\section{Template L2 di riferimento}

\subsection{MacroRequirement}
\begin{lstlisting}
# <MR-ID> - <Title>
Lifecycle: <...>
Authority: <...>   [if used by corpus governance]

## Intent
...
## Context
...
## Stakeholders
...
## Scope
IN: ...
OUT: ...
## Assumptions / Constraints
...
## dependsOn
<MR-ID or None>
\end{lstlisting}

\subsection{Decision}
\begin{lstlisting}
# <Decision-ID> - <Title>
Parent MacroRequirement: <MR-ID>
Decision kind: <SELECTABLE / NECESSITY-CONSTRAINED / DEFAULT>
              [authoring/review classification]

## Context
...
## Decision
...
## Consequences
...
\end{lstlisting}

\subsection{FunctionalRequirement}
\begin{lstlisting}
# <FR-ID> - <Title>
Parent Decision: <Decision-ID>

## Functional obligation
...
## Normative clauses
NC-1: <Subject> MUST ...
...
## SPO references
<Subject> -- <Predicate> --> <Object>
\end{lstlisting}

\subsection{SpecializedRequirement / SecurityRequirement}
\begin{lstlisting}
# <SR-or-SEC-ID> - <Title>
Parent FunctionalRequirement: <FR-ID>
protectedSecurityProperty: <Property>   [SecurityRequirement only]

## Normative clauses
...

# Optional L2 review aid, not second source of truth:
Failure mode: ...
\end{lstlisting}

\section{R43 non-normative conclusion}

\begin{ddtacore}[title={Cosa consolida questa candidate}]
\begin{lstlisting}
R5 authoring principles = preserved
native DDTA authoring = preserved
L1/L2/L3/L4 boundary = preserved
source-first / authority-first discipline = preserved
technical information preservation = strengthened by R25 evidence
condition / selection / pipeline lessons = explicit
DermaTriage source gaps = preserved as source gaps
downstream BA candidates/open pressures = boundary-only, not authoring constructs
R5 = current authority until governed promotion
\end{lstlisting}
\end{ddtacore}

\end{document}
